\newif\ifglobal

%%%%%%%%%%%%%%%%%%%%%%%
%%% Compile to view %%%
%%%%%%%%%%%%%%%%%%%%%%%

%%%%%%%%%%%%%%%%%%%%%%%%%%%%%%%%%%%%%%%%%%%%%%%%%%%
%%% Readme:                                     %%%
%%% https://www.overleaf.com/read/bwdjvfjksvjy/ %%%
%%%%%%%%%%%%%%%%%%%%%%%%%%%%%%%%%%%%%%%%%%%%%%%%%%%

% >>> M?: marks a module setting number or important notice

% >>> M4: Set {./relative-path} to external files for this module <<<
\def \modulefiles {./35-RNG}
% To add an external file, use:
% (Replace '---' with actual filename)
% '\input{\modulefiles/tables/---.tex}' for tables
% '\includegraphics{\modulefiles/figures/---}' for figures
% '\subfile{\modulefiles/---}' for register files


% >>> M1: This module is in global project? <<<
% 'YES' - uncomment; 'NO' - comment out
% \globaltrue

\ifglobal
    \documentclass[main\_\_CN.tex]{subfiles}

\else
    % Font "FandolSong-Regular" used by ctex does not have GSUB table.
    % It causes warnings that do not affect compilation and can be ignored.
    \PassOptionsToPackage{quiet}{fontspec}

    \documentclass[openany, 10pt]{book}

    % >>> M2: In ./chip-spec-settings/preamble-trm-module, also find
    %         settings for visibility of todo notes, watermark <<<
    \usepackage{preamble-trm-module}


    % Enable Chinese typesetting
    \usepackage{ctex}

    % >>> M3: Temporary labels in Chinese
    %         you can add your own labels there <<<
    \myexternaldocument{temp-labels-cn}

\fi %\ifglobal


\setcounter{secnumdepth}{4}



% >>> M5: If this module needs any updates in
% './chip-spec-settings/preamble-trm-module.sty'
% (include additional package, etc.), contact your module PM
% or create the file './chip-spec-settings/changelog.tex' make the
% necessary updates yourself and document them in this file <<<

\raggedright


\begin{document}


% Sets language into which pre-defined bits of text are translated
\selectlanguage{chinese}

% Do not delete this block
\ifglobal
% Add the following front matter if
% this module is outside of global project
\else
    \InputIfFileExists{readme}{}{}
    {\let\clearpage\relax \listoftodos}
    {\let\clearpage\relax \listofdonetodos}
    \tableofcontents
    %\listoftables
    %\listoffigures
\fi


%%%%%%%%%%%%%%%%%%%%%%%%%
%%% TRM Module Begins %%%
%%%%%%%%%%%%%%%%%%%%%%%%%

\hypertarget{rng}{}
\chapter{随机数发生器 (RNG)}\label{mod:rng}

\section{概述}\label{sec:rng-overview}

\chipname{} 内置真随机数发生器，可产生两类随机数：CRC 随机数和 NIST 随机数。两类随机数均可作为加密操作的基础并且均通过了 dieharder 测试。

%%%%%%%%%%%%%%%%%%%
%%% Terminology %%%
%%%%%%%%%%%%%%%%%%%

\section{术语}
\label{sec:rng-term}

为了更好地说明随机数发生器 (RNG) 的功能，本章使用了以下术语。

\begin{longtable}[c]{ R{4cm} p{12cm} }
    \textbf{CRC 随机数} & 真随机数发生器产生的随机数，软件可从寄存器 \hyperref[regdesc:TRNGCRCSYNCDATAREG]{TRNG\_CRC\_SYNC\_DATA\_REG} 读取。\\
    \textbf{NIST 随机数} & 真随机数发生器产生的随机数，供芯片安全子系统内的加密模块使用。该类随机数由符合 NIST SP 800-90B 的片上健康测试电路检测。\\
    \textbf{NIST SP 800-90B} & 美国国家标准与技术研究院 (National Institute of Standards and Technology) 特别出版物 800-90B。该出版物定义了 RNG 使用的健康测试：Startup Test、Repetition Test 和 Adaptive Test。\\
\end{longtable}

\section{主要特性}
\label{sec:rng-features}

\begin{itemize}
    \item 热噪声、异步时钟和环形振荡器作为噪声源
    \item 软件可读的 CRC 随机数
    \item 面向芯片安全子系统内加密模块的 NIST 随机数
    \item 符合 NIST SP 800-90B 的片上健康测试电路（具体测试请参考 \ref
    {sec:rng-nist-rng} \nameref{sec:rng-nist-rng} 小节）
\end{itemize}

\section{功能描述}\label{sec:rng-func-descr}

CRC 随机数和 NIST 随机数都可使用\textbf{热噪声}、\textbf{异步时钟}和\textbf{环形振荡器}作为噪声输入。此外，NIST 通路还支持复用 CRC\_DATA 作为输入。

\begin{itemize}
    \item \textbf{热噪声}可以来自 SAR\_ADC 或 RF\_ADC。当 SAR\_ADC 工作时，SAR\_ADC 产生比特流，并通过异或 (XOR) 逻辑运算作为随机数种子进入随机数发生器。
    \item RC\_FAST\_CLK 是一种\textbf{异步时钟源}，会产生电路亚稳态。这种亚稳态也可以作为随机数种子，进入随机数发生器。关于 RC\_FAST\_CLK，详见以下 \hyperref[others:rng-note-clock-seed]{说明}。
    \item BUF\_CELL（缓存单元）实现的\textbf{环形振荡器} (BUF\_CHAIN) 生成的计数值作为随机数种子进入随机数发生器。
\end{itemize}
噪声源结构如图 \ref{fig:rng-noise-source} \textit{\nameref{fig:rng-noise-source}} 所示。
\vspace{2em}
\begin{figure}[H]
    \centering
    \includegraphics[width=1.0\textwidth]{\modulefiles/figures/35-RNG-noise-source.png}
    \caption{噪声源}
    \label{fig:rng-noise-source}
\end{figure}

软件可以从寄存器 \hyperref[regdesc:TRNGCRCSYNCDATAREG]{TRNG\_CRC\_SYNC\_DATA\_REG} 中读取 CRC 随机数。NIST 随机数的质量高于 CRC 随机数。NIST 随机数仅能通过芯片安全子系统内的加密模块（例如 Key Manager）请求获取，不能通过软件通路直接读取。

当 SAR\_ADC 打开时，每个 RC\_FAST\_CLK 时钟周期内（来自内部 RC 振荡器，请参见章节 \ref{mod:resclk} \textit{\nameref{mod:resclk}}），随机数发生器将获得 1 位熵。因此，为了获得最大熵值，乐鑫建议读取 \hyperref[regdesc:TRNGCRCSYNCDATAREG]{TRNG\_CRC\_SYNC\_DATA\_REG} 寄存器的速率不超过 1 MHz。更高速率读取可能会降低随机性质量。


\section{配置流程}\label{sec:rng-programm-descr}


\subsection{CRC RNG}

在使用 \chipname{} 的随机数发生器时，在 SoC 层面应确保在任何 TRNG 操作之前完成以下步骤：

\begin{enumerate}
    \item 配置寄存器 \hyperref[regdesc:LPPERICLKRSTRNGCTRLREG]{LP\_PERICLKRST\_RNG\_CTRL\_REG} 中的 \hyperref[fielddesc:LPPERICLKRSTLPRNGCKEN]{LP\_PERICLKRST\_LP\_RNG\_CLK\_EN} 为 1，以打开 RNG 时钟和用作熵源的 RC\_FAST\_CLK。
    \item 配置 \hyperref[fielddesc:LPPERICLKRSTLPRNGRSTEN]{LP\_PERICLKRST\_LP\_RNG\_RST\_EN} 为 1，再配置为 0，以产生 RNG IP 的复位脉冲。
\end{enumerate}

随后使能 BUF\_CHAIN。否则可能产生伪随机数。也可将 SAR\_ADC 和 RC\_FAST\_CLK 作为额外噪声源使能。
\begin{itemize}
    \item 配置寄存器 \hyperref[regdesc:TRNGCONFREG]{TRNG\_CONF\_REG} 中的 \hyperref[fielddesc:TRNGSAMPLEENABLE]{TRNG\_SAMPLE\_ENABLE} 使能 BUF\_CHAIN。
    \item 关于 SAR\_ADC，请参见章节 \ref{mod:adc} \textit{\nameref{mod:adc}}。
    \item 配置寄存器 \hyperref[regdesc:PMUHPSLEEPLPCKPOWERREG]{PMU\_HP\_SLEEP\_LP\_CK\_POWER\_REG} 中的 \hyperref[fielddesc:PMUHPSLEEPXPDFOSCCLK]{PMU\_HP\_SLEEP\_XPD\_FOSC\_CLK} 使能 RC\_FAST\_CLK。关于 RC\_FAST\_CLK，详见以下 \hyperref[others:rng-note-clock-seed]{说明}。
\end{itemize}

随后将 \hyperref[fielddesc:TRNGNOISECRCEN]{TRNG\_NOISE\_CRC\_EN} 置 1，以启用 CRC 随机数生成。

\begin{tiplisting}
    \label{others:rng-note-clock-seed}
    RC\_FAST\_CLK 可以提高 RNG 熵值。然而，为了保证足够熵值，在将 RC\_FAST\_CLK 用作噪声源时，乐鑫仍建议保持 BUF\_CHAIN 使能。
\end{tiplisting}

在软件使用随机数发生器时，请多次读取 \hyperref[regdesc:TRNGCRCSYNCDATAREG]{TRNG\_CRC\_SYNC\_DATA\_REG} 寄存器的值，直至获得足够多的随机数。在读取此寄存器时，注意控制速率，不要超过第 \ref{sec:rng-func-descr} 节 \textit{\nameref{sec:rng-func-descr}} 所述限制。

\subsection{NIST RNG}\label{sec:rng-nist-rng}

对于 NIST 随机数，其噪声源支持通过 \hyperref[regdesc:TRNGCONFREG]{TRNG\_CONF\_REG} 寄存器中的 \hyperref[fielddesc:TRNGNOISESOURCESEL]{TRNG\_NOISE\_SOURCE\_SEL} 以及对应的 \hyperref[fielddesc:TRNGNOISEPOSSEL]{TRNG\_NOISE\_POS\_SEL} 进行选择：
\begin{itemize}
    \item 复用 CRC 随机数作为 NIST 健康测试的输入随机数据：配置 \hyperref[fielddesc:TRNGNOISESOURCESEL]{TRNG\_NOISE\_SOURCE\_SEL}[4] = 1。若选择此配置，则会忽略下列所有噪声源。
    \item 使能 SAR\_ADC：配置 \hyperref[fielddesc:TRNGNOISESOURCESEL]{TRNG\_NOISE\_SOURCE\_SEL}[3] = 1。
    \item 使能 RF\_ADC：配置 \hyperref[fielddesc:TRNGNOISESOURCESEL]{TRNG\_NOISE\_SOURCE\_SEL}[2] = 1。
    \item 使能 RC\_FAST\_CLK：配置 \hyperref[fielddesc:TRNGNOISESOURCESEL]{TRNG\_NOISE\_SOURCE\_SEL}[1] = 1。
    \item 使能 BUF\_CHAIN：配置 \hyperref[fielddesc:TRNGNOISESOURCESEL]{TRNG\_NOISE\_SOURCE\_SEL}[0] = 1。
\end{itemize}
除 CRC\_DATA 通路外，上述噪声源支持同时选择多个输入通道。为保证随机数输出质量，若只选择噪声源 BUF\_CHAIN 或 RC\_FAST\_CLK，\hyperref[fielddesc:TRNGHEALTHTESTERRORCODE]{TRNG\_HEALTH\_TEST\_ERROR\_CODE} 会记录错误代码 1 并且自行清除当前 RNG 模块内存放的 NIST 随机数来防止输出不合规随机数。出错后建议软件将 \hyperref[fielddesc:LPPERICLKRSTLPRNGRSTEN]{LP\_PERICLKRST\_LP\_RNG\_RST\_EN} 置 1 清除 RNG 当前所有状态然后重新运行一遍启动 NIST RNG 通路的步骤。

对于 \hyperref[fielddesc:TRNGNOISEPOSSEL]{TRNG\_NOISE\_POS\_SEL}，该字段采用独热编码。不同于 \hyperref[fielddesc:TRNGNOISESOURCESEL]{TRNG\_NOISE\_SOURCE\_SEL}，该字段只能置位其中一位。应确保 \hyperref[fielddesc:TRNGNOISEPOSSEL]{TRNG\_NOISE\_POS\_SEL} 与 \hyperref[fielddesc:TRNGNOISESOURCESEL]{TRNG\_NOISE\_SOURCE\_SEL} 中已使能的某一路噪声源匹配，否则 RNG 无法正确输出随机数。

为保证 NIST 随机数输出符合 NIST SP 800-90B 标准，RNG 中内置了三套 NIST SP 800-90B 标准中定义的强制测试：Startup Test、Repetition Test 和 Adaptive Test。对于 \hyperref[fielddesc:TRNGREPETITIONVALUEC]{TRNG\_REPETITION\_VALUE\_C} 和 \hyperref[fielddesc:TRNGADAPTIVEVALUEC]{TRNG\_ADAPTIVE\_VALUE\_C} 的取值，请根据随机数长度和位宽，参考 NIST SP 800-90B 手册中的建议值。
\begin{itemize}
    \item Startup Test：RNG 启动后检测一段时间采样的随机数熵源输出数据，防止出现启动后熵源不稳定导致输出随机数失真的情况。配置字段 \hyperref[fielddesc:TRNGSTARTUPTESTSTART]{TRNG\_STARTUP\_TEST\_START} 为 1 开始此测试。测试的数据长度为 512 + \hyperref[fielddesc:TRNGSTARTUPTESTLIMIT]{TRNG\_STARTUP\_TEST\_LIMIT} 个数据。
    \item Repetition Test：检测输出的随机数是否被卡在了某一固定值，通过计数连续出现多少相同的随机数输出来判断随机数质量是否达标。此测试的阈值可以通过配置字段 \hyperref[fielddesc:TRNGREPETITIONVALUEC]{TRNG\_REPETITION\_VALUE\_C} 进行设置，最大支持配置 511。配置值越低测试越严格。
    \item Adaptive Test：检测在任意连续的 512 个生成的随机数窗口内是否存在某一随机数出现个数过高的情况。出现个数限制可以通过 \hyperref[fielddesc:TRNGADAPTIVEVALUEC]{TRNG\_ADAPTIVE\_VALUE\_C} 设置。
\end{itemize}

仅在测试模式下可将 \hyperref[fielddesc:TRNGHEALTHTESTBYPASS]{TRNG\_HEALTH\_TEST\_BYPASS} 置 1 以旁路上述健康测试。其余情况应始终将 \hyperref[fielddesc:TRNGHEALTHTESTBYPASS]{TRNG\_HEALTH\_TEST\_BYPASS} 置 0。
\\
若软件希望读取 NIST 随机数则需要在配置完 NIST 随机数相关配置且开启 Startup Test 后配置寄存器 \hyperref[fielddesc:TRNGSWRANDOMREQ]{TRNG\_SW\_RANDOM\_REQ} 然后等待寄存器 \hyperref[fielddesc:TRNGSWDATAVALID]{TRNG\_SW\_DATA\_VALID} 为 1 后读取寄存器 \hyperref[fielddesc:TRNGSWRANDOMDATA]{TRNG\_SW\_RANDOM\_DATA} 获取 NIST 随机数。



% >>> If your module has no registers, delete sample text
%       for Sections Register Summary and Registers below <<<

%%%%%%%%%%%%%%%%%%%%%%%%
%%% Register Summary %%%
%%%%%%%%%%%%%%%%%%%%%%%%

% >>> Update the sample text below and insert
%     the info on register summary into the subfile <<<

\newpage

\hypertarget{rng-reg-summ}{}
\section{寄存器列表}
\label{sec:rng-reg-summ}

请查看章节 \hyperref[glossary-access-types]{\textit{寄存器的访问类型}}，了解“访问”列缩写的含义。

\subfile{\modulefiles/35-RNG-reg-summ__CN}



%%%%%%%%%%%%%%%%%
%%% Registers %%%
%%%%%%%%%%%%%%%%%

% >>> Update the sample text below and insert
%      the info on registers into the subfile <<<

\newpage
\section{寄存器}
\label{sec:rng-reg}

关于保留 (reserved) 域的处理，请查看章节 \hyperref[programming-reserved-field]{如何配置寄存器的保留域}。

\subfile{\modulefiles/35-RNG-reg__CN}

\end{document}
